\chapter{\MakeUppercase{Hasil Pengujian dan Analisis}}

Bab ini menyajikan hasil implementasi, pengujian, serta analisis terhadap sistem yang telah dirancang. Tujuan pengujian adalah untuk memverifikasi bahwa sistem telah bekerja sesuai dengan spesifikasi, kebutuhan, dan tujuan penelitian yang telah ditetapkan.

Pada bab ini, hasil pengujian tidak hanya ditampilkan dalam bentuk data, tabel, grafik, atau gambar, tetapi juga harus dianalisis secara kritis untuk menjelaskan makna dari hasil yang diperoleh. Analisis yang dilakukan harus mampu menunjukkan tingkat keberhasilan sistem dalam menyelesaikan permasalahan yang dirumuskan pada penelitian.

Secara umum, isi bab ini dapat mencakup beberapa bagian berikut:

\begin{enumerate}[topsep=0pt,itemsep=0pt,parsep=0pt,leftmargin=*]
\item Deskripsi implementasi sistem yang telah direalisasikan berdasarkan rancangan pada bab sebelumnya.
\item Metode dan skenario pengujian yang digunakan.
\item Hasil pengujian setiap bagian atau fungsi sistem.
\item Hasil pengujian sistem secara keseluruhan.
\item Analisis terhadap hasil pengujian.
\item Perbandingan hasil yang diperoleh dengan target, spesifikasi, teori, penelitian terdahulu, atau standar yang digunakan sebagai acuan.
\item Identifikasi kelebihan, keterbatasan, dan faktor-faktor yang memengaruhi kinerja sistem.
\end{enumerate}

Dalam melakukan analisis, beberapa pertanyaan berikut sebaiknya dapat dijawab:

\begin{itemize}[topsep=0pt,itemsep=0pt,parsep=0pt,leftmargin=*]
\item Apakah sistem telah bekerja sesuai dengan rancangan?
\item Apakah tujuan penelitian telah tercapai?
\item Apakah hasil yang diperoleh sesuai dengan teori atau referensi yang digunakan?
\item Faktor apa saja yang memengaruhi hasil pengujian?
\item Apa kelebihan dan keterbatasan sistem yang dikembangkan?
\item Bagaimana hasil penelitian dibandingkan dengan metode, sistem, atau penelitian lain yang sejenis?
\end{itemize}

Setiap hasil pengujian yang disajikan harus didukung oleh data yang memadai. Data dapat disajikan dalam bentuk tabel, grafik, gambar, tangkapan layar, maupun bentuk lain yang sesuai. Seluruh data yang ditampilkan harus dirujuk dan dijelaskan dalam uraian sehingga pembaca dapat memahami konteks serta makna dari hasil yang diperoleh.

Hal-hal yang perlu diperhatikan dalam penyusunan bab ini adalah sebagai berikut:

\begin{itemize}[topsep=0pt,itemsep=0pt,parsep=0pt,leftmargin=*]
\item Jangan hanya menampilkan data tanpa analisis.
\item Analisis harus didasarkan pada fakta dan data yang diperoleh, bukan asumsi semata.
\item Setiap tabel, grafik, gambar, atau hasil pengujian harus dirujuk dan dijelaskan.
\item Hasil yang tidak sesuai harapan tetap harus dilaporkan dan dianalisis penyebabnya.
\item Analisis harus mengarah pada pembuktian bahwa tujuan penelitian telah atau belum tercapai.
\end{itemize}

Pada akhir bab, pembaca harus dapat memahami tingkat keberhasilan sistem yang dikembangkan, faktor-faktor yang memengaruhi kinerjanya, serta kontribusi hasil penelitian terhadap penyelesaian permasalahan yang diteliti.
